% =====================================================================
%  A Comparative Taxonomy of Market Regimes
%  Seminar paper — Chair of Financial Management and Capital Markets
%  Compile with:  pdflatex paper.tex  (run TWICE for TOC / refs / citations)
%  Figures are read from the results/ directory.
% =====================================================================
\documentclass[12pt,a4paper]{article}

\usepackage[utf8]{inputenc}
\usepackage[T1]{fontenc}
\usepackage{amsmath,amssymb}
\usepackage{newtxtext,newtxmath}      % Times New Roman–like text and math
\usepackage[left=2.5cm,right=2.5cm,top=2.5cm,bottom=2cm]{geometry}
\usepackage{setspace}
\usepackage{graphicx}
\usepackage{booktabs}
\usepackage{array}
\usepackage{caption}
\usepackage{microtype}
\usepackage{float}
\usepackage{fancyhdr}
\usepackage[authoryear,round]{natbib}
\usepackage[hidelinks]{hyperref}

\graphicspath{{results/figures_light/}{results/}{./results/}}
\captionsetup{font=small,labelfont=bf}
\onehalfspacing

% Page numbers centered at the bottom of every page (cover page stays unnumbered)
\pagestyle{fancy}
\fancyhf{}
\fancyfoot[C]{\thepage}
\renewcommand{\headrulewidth}{0pt}
\fancypagestyle{plain}{%
  \fancyhf{}%
  \fancyfoot[C]{\thepage}%
  \renewcommand{\headrulewidth}{0pt}%
}
\setlength{\parindent}{0pt}
\setlength{\parskip}{0.6em}

% Small helper for the "Notes / Source" block under floats
\newcommand{\floatnote}[1]{\par\vspace{2pt}{\footnotesize\textit{#1}}}
\newcommand{\dvix}{\Delta\mathrm{VIX}}
\newcommand{\rf}{r^{f}}

\begin{document}

% =====================================================================
%  COVER PAGE
% =====================================================================
\begin{titlepage}
\centering
{\large [University Name]\par}
{\large Chair of Financial Management and Capital Markets\par}
\vspace{1.2cm}
{\large Seminar Paper\par}
\vspace{2.2cm}
{\Large\bfseries A Comparative Taxonomy of Market Regimes:\\[0.3em]
Machine Learning versus Hidden Markov Models in\\[0.3em]
US Tactical Sector Allocation\par}
\vspace{2.2cm}
\begin{tabular}{rl}
Submitted by:        & [Your Full Name]\\
Matriculation No.:   & [XXXXXXX]\\
Study Programme:     & [B.Sc./M.Sc. \ldots]\\
Supervisor:          & [Supervisor Name]\\
Semester:            & [e.g.\ Summer Semester 2026]\\
\end{tabular}
\vfill
{\normalsize [Place], \today\par}
\end{titlepage}

% =====================================================================
%  FRONT MATTER (roman page numbers)
% =====================================================================
\pagenumbering{roman}

\tableofcontents
\clearpage

\addcontentsline{toc}{section}{List of Figures}
\listoffigures
\clearpage

\addcontentsline{toc}{section}{List of Tables}
\listoftables
\clearpage

% ---- List of Abbreviations ------------------------------------------
\addcontentsline{toc}{section}{List of Abbreviations}
\section*{List of Abbreviations}
\begin{tabular}{ll}
AIC   & Akaike Information Criterion \\
BIC   & Bayesian Information Criterion \\
CAPM  & Capital Asset Pricing Model \\
DSR   & Deflated Sharpe Ratio \\
EM    & Expectation--Maximisation \\
ETF   & Exchange-Traded Fund \\
FRED  & Federal Reserve Economic Database \\
GMM   & Gaussian Mixture Model \\
HMM   & Hidden Markov Model \\
MVP   & Minimum-Variance Portfolio \\
OLS   & Ordinary Least Squares \\
PSR   & Probabilistic Sharpe Ratio \\
SLSQP & Sequential Least-Squares Quadratic Programming \\
SPDR  & Standard \& Poor's Depositary Receipts \\
TPF   & Tangency Portfolio (maximum-Sharpe) \\
VIX   & CBOE Volatility Index \\
\end{tabular}
\clearpage

% ---- List of Symbols ------------------------------------------------
\addcontentsline{toc}{section}{List of Symbols}
\section*{List of Symbols}
\begin{tabular}{ll}
$P_{i,t}$            & Adjusted closing price of asset $i$ at time $t$ \\
$r_{i,t}$            & Log return of asset $i$ \\
$e_{i,t}$            & Excess (over risk-free) return of asset $i$ \\
$\rf_t$              & Daily risk-free rate \\
$\dvix_t$           & Daily change in the VIX index \\
$s_t$               & Latent regime (state) at time $t$ \\
$\mu_k,\ \Sigma_k$  & State-$k$ conditional mean vector and covariance matrix \\
$A_{ij}$            & Markov transition probability from state $i$ to state $j$ \\
$\alpha,\ \beta$    & CAPM intercept (Jensen's alpha) and market beta \\
$\varepsilon_{i,t}$ & Regression residual \\
$w$                 & Vector of portfolio weights \\
$\tau_t$            & One-way portfolio turnover at rebalance $t$ \\
$V,\ p$             & Reality-Check statistic and bootstrap $p$-value \\
$T,\ N$             & Number of time observations and of assets \\
\end{tabular}
\clearpage

% =====================================================================
%  MAIN TEXT (arabic page numbers)
% =====================================================================
\pagenumbering{arabic}

\begin{abstract}
\noindent
We compare three families of market-regime classifiers---a rule-based volatility
trigger, a Gaussian Hidden Markov Model (HMM), and an unsupervised Gaussian Mixture
Model (GMM)---all feeding an identical regime-conditional portfolio construction, and
ask whether conditioning a US sector-rotation strategy on the
inferred regime improves out-of-sample, after-cost performance relative to a passive
S\&P~500 benchmark. Using daily data on nine SPDR sector ETFs from 2004 to 2025, all
regime models are estimated in a strict expanding-window walk-forward design, and
portfolios are formed each month by regime-conditional minimum-variance (MVP) and
maximum-Sharpe tangency (TPF) optimisation with Ledoit--Wolf covariance shrinkage. We
find that regime conditioning delivers a large and robust reduction in volatility and
drawdown---the minimum-variance strategies cut the maximum drawdown from roughly $-60\%$
(S\&P~500) to about $-37\%$ and raise the Sharpe ratio from $0.44$ to $0.48$--$0.54$---but
that this gain comes entirely from risk targeting. Because all three classifiers feed an
identical portfolio construction, the tangency (return-timing) variants fail uniformly, and
no strategy out-returns the benchmark once data snooping is corrected for (White's Reality
Check $p=0.98$); the best strategy's risk-adjusted edge is only marginally significant after
deflation and comes entirely from lower risk, not higher return. We conclude that the
economic value of regime-based tactical allocation in this
setting is concentrated in risk control rather than abnormal returns, and that it is largely
insensitive to the sophistication of the regime classifier.
\end{abstract}

% =====================================================================
\section{Introduction}
% =====================================================================
A long tradition in empirical finance holds that asset returns are generated by a small
number of recurring ``states of the world''---calm bull markets, turbulent transitions,
and crises---with sharply different first and second moments \citep{hamilton1989,
ang2002, guidolin2007}. If such regimes are identifiable in real time, an investor could
in principle rotate across assets to harvest the favourable states and de-risk in the
unfavourable ones. This paper studies that proposition for US equity \emph{sectors}, the
most liquid vehicle for expressing tactical macro views, and confronts a practical
question that much of the academic literature elides: once a regime model is estimated
only on information available at the time of trading, and once realistic transaction
costs are charged, does regime conditioning actually beat a passive benchmark?

\paragraph{Objective and research question.}
The objective of this paper is to quantify, on a strictly out-of-sample and after-cost
basis, the economic value added by market-regime conditioning in US sector allocation,
and to decompose that value into a \emph{risk} component and a \emph{return} component.
We address three research questions. (RQ1) Does a data-driven latent-state model (HMM or
machine learning) identify regimes more usefully than a transparent VIX rule? (RQ2) Is
any incremental value attributable to regime \emph{persistence}, isolated by contrasting
the Markovian HMM with a memoryless mixture? (RQ3) Does the best-performing strategy
survive a formal data-snooping correction, or is its outperformance a selection artefact?

\paragraph{Contributions.}
We make three contributions. First, we place three structurally distinct regime
classifiers on a common, strictly out-of-sample footing: a transparent VIX threshold
rule, a Gaussian HMM that explicitly models regime \emph{persistence}, and an unsupervised
Gaussian Mixture Model (GMM) that clusters days into regimes without any temporal
dependence. All three feed an \emph{identical} downstream portfolio construction---moments
estimated from the historical days in the prevailing regime---so that the only thing that
differs is the regime-labelling method, and performance gaps are cleanly attributable to
it. The VIX-versus-model contrast isolates the value of a data-driven latent-state model
over a fixed rule, while the HMM-versus-GMM contrast isolates the value of modelling
temporal dependence (the HMM's Markov transition matrix) versus a memoryless mixture.
Second, we evaluate each regime signal through two portfolio lenses---minimum-variance
(which depends only on the covariance matrix) and maximum-Sharpe tangency (which
additionally depends on expected returns)---so as to separate the value of risk estimation
from the value of return estimation. Third, and most importantly for the credibility of the
exercise, we measure portfolio turnover directly, charge transaction costs, and subject
the leading strategy to formal multiple-testing corrections, because evaluating many
strategies on one sample mechanically inflates the best in-sample result
\citep{harvey2016}.

Our headline finding is a clean dissociation. Regime conditioning reliably \emph{reduces
risk}: across all three classifiers, the minimum-variance variants achieve Sharpe ratios
of $0.46$--$0.50$ against the S\&P~500's $0.44$ at roughly two-thirds of the volatility
and three-fifths of the drawdown, with low turnover that is robust to plausible trading
costs. But the one strategy that appears to \emph{add return}---the machine-learning
tangency portfolio, with a $9.5\%$ annualised return and a $4.3\%$ CAPM alpha---trades
more than six times its capital per year, and its return advantage over the S\&P~500 is
statistically indistinguishable from chance once data snooping is accounted for. The
value of ``intelligence'' in regime modelling, in this dataset, is defensive.

\paragraph{Outline.}
The remainder proceeds as follows. Section~\ref{sec:lit} reviews the relevant literature
on regime-switching models, volatility-based classification, machine learning for
returns, and performance evaluation. Section~\ref{sec:data} describes the data.
Section~\ref{sec:method} develops the methodology---feature construction, the three
regime classifiers, the conditional CAPM, portfolio construction, and the evaluation
framework. Section~\ref{sec:results} presents the empirical results. Section~\ref{sec:disc}
discusses their interpretation, Section~\ref{sec:limit} states limitations, and
Section~\ref{sec:concl} concludes.

% =====================================================================
\section{Related Literature}\label{sec:lit}
% =====================================================================
\paragraph{Regime-switching models.}
The Markov-switching framework of \citet{hamilton1989} formalised the idea that financial
time series alternate between discrete states with state-dependent dynamics.
\citet{ang2002} and \citet{guidolin2007} applied multivariate regime-switching models to
equity allocation and documented that correlations and volatilities rise sharply in
``bear'' regimes, eroding diversification precisely when it is most needed. The Gaussian
HMM we employ is the canonical estimation device for such latent-state models, fitted by
the Baum--Welch expectation--maximisation algorithm \citep{baum1970, dempster1977} and
decoded by the \citet{viterbi1967} algorithm.

\paragraph{Volatility as a regime proxy.}
The CBOE VIX index, the risk-neutral expected volatility of the S\&P~500 over the next 30
days \citep{whaley2009}, is the most widely used real-time stress gauge. Practitioner and
academic conventions treat VIX below 20 as calm, 20--30 as elevated, and above 30 as
crisis \citep{whaley2009, bloom2009}, a mapping we adopt as our transparent benchmark
classifier. \citet{bloom2009} shows that volatility shocks are themselves a driver of
real and financial dynamics, motivating the use of the \emph{change} in VIX as a feature.

\paragraph{Unsupervised clustering of regimes.}
Gaussian Mixture Models provide a soft-clustering analogue to the HMM without the Markov
persistence assumption \citep{mclachlan2000}. Fitting a finite Gaussian mixture to the
joint distribution of volatility changes and sector returns partitions the sample into
latent states in the same emission family as the HMM, but treats each observation
independently. Comparing the two therefore isolates the contribution of the temporal
(Markov) structure while holding the distributional assumptions fixed.

\paragraph{Portfolio construction and estimation risk.}
\citet{markowitz1952} mean--variance optimisation is notoriously sensitive to estimation
error in expected returns, to the point that the na\"ive $1/N$ portfolio is hard to beat
out of sample \citep{demiguel2009}. Two responses inform our design: minimum-variance
investing, which avoids the expected-return input entirely and has delivered strong
risk-adjusted performance empirically \citep{clarke2006}, and covariance shrinkage, for
which \citet{ledoit2004} derive an optimal linear combination of the sample covariance
and a structured target.

\paragraph{Performance evaluation under data snooping.}
Because we search over many strategies, in-sample outperformance is upward biased.
\citet{white2000} introduced the Reality Check, a bootstrap test of whether the
\emph{best} of many strategies beats a benchmark, implemented here with the stationary
bootstrap of \citet{politis1994}. \citet{bailey2014} propose the Deflated Sharpe Ratio,
which benchmarks a strategy's Sharpe against the expected maximum Sharpe attainable by
chance across the trials searched. \citet{harvey2016} make the general case for raising
significance thresholds in light of the multiplicity of tested strategies.

% =====================================================================
\section{Data}\label{sec:data}
% =====================================================================
The sample comprises daily data from January 2004 to December 2025 (5{,}533 trading
days), sourced from Yahoo!\ Finance and the Federal Reserve Economic Database (FRED) and
cached locally for reproducibility.

\paragraph{Asset universe.}
The nine SPDR Select Sector ETFs---Technology (XLK), Consumer Discretionary (XLY),
Financials (XLF), Industrials (XLI), Materials (XLB), Energy (XLE), Consumer Staples
(XLP), Health Care (XLV), and Utilities (XLU)---span the investable US equity market,
with SPDR S\&P~500 (SPY) as the benchmark. We use dividend- and split-adjusted closing
prices.

\paragraph{State and macro variables.}
The CBOE VIX index proxies market stress. The risk-free rate is the 13-week
Treasury-bill annualised yield (\texttt{\textasciicircum IRX}), converted to a daily
decimal rate. Two macro--financial conditioning variables are drawn from FRED: the TED
spread (\texttt{TEDRATE}), a funding-stress indicator, and the 10-year--2-year Treasury
term spread (\texttt{T10Y2Y}), a leading indicator of the business cycle. Log returns and
excess returns are defined as
\begin{equation}
  r_{i,t}=\ln\!\left(\frac{P_{i,t}}{P_{i,t-1}}\right),
  \qquad
  e_{i,t}=r_{i,t}-\rf_t,
\end{equation}
where $e_{i,t}$ is the excess return over the risk-free rate.

% =====================================================================
\section{Methodology}\label{sec:method}
% =====================================================================

\subsection{Feature construction}
Both data-driven classifiers operate on the same regime-feature vector: the daily change
in VIX, $\dvix_t=\mathrm{VIX}_t-\mathrm{VIX}_{t-1}$, stacked with the nine sector log
returns (a 10-dimensional vector). Volatility---through the level and the change of the
VIX---is the primary driver of which regime prevails, so a parsimonious common feature set
keeps the HMM and the GMM comparable and guards against overfitting the state assignment.

\subsection{Regime identification}

\paragraph{(a) VIX rule.}
The transparent benchmark assigns each day to Calm ($\mathrm{VIX}<20$), Transitional
($20\le\mathrm{VIX}<30$), or Crisis ($\mathrm{VIX}\ge30$), following \citet{whaley2009}
and \citet{bloom2009}.

\paragraph{(b) Hidden Markov Model.}
We fit a three-state Gaussian HMM in which the observed feature vector is, conditional on
the latent state $s_t\in\{0,1,2\}$, multivariate normal,
\begin{equation}
  x_t \mid s_t=k \;\sim\; \mathcal{N}(\mu_k,\Sigma_k),
  \qquad
  A_{ij}=\Pr(s_{t+1}=j\mid s_t=i),
\end{equation}
with full state-specific covariances $\Sigma_k$ and a first-order Markov transition
matrix $A$. Parameters are estimated by Baum--Welch EM; because EM converges only to
local optima, we take the best of twenty random restarts by log-likelihood, and recover
the most probable state path by Viterbi decoding. HMM state indices are not identified by
the likelihood (the ``label-switching'' problem), so we sort the states by their
state-conditional mean $\dvix$, fixing state $0$ as Calm and state $2$ as Crisis and
permuting $A$, $\mu_k$ and $\Sigma_k$ accordingly.

\paragraph{(c) Gaussian Mixture Model.}
The machine-learning classifier is a three-component Gaussian Mixture fitted on the regime
features, which assigns each day a hard cluster label sorted by $\dvix$ for
identifiability. Unlike the HMM, the mixture is \emph{memoryless}: it has no transition
matrix and treats each day independently. This is deliberate---contrasting the GMM with
the HMM isolates the marginal value of modelling regime \emph{persistence}, holding the
Gaussian emission structure fixed. Crucially, the GMM produces regime \emph{labels only};
the portfolio moments it feeds are estimated, exactly as for the VIX rule and the HMM,
from the historical days assigned to the prevailing regime (Section~\ref{sec:method}). No
auxiliary return-forecasting model is interposed, so that any performance difference
across classifiers is attributable to the regime-labelling method itself rather than to a
separate return predictor.

\paragraph{(d) Walk-forward estimation.}
Both learned models are fitted in an \emph{expanding-window walk-forward}: trained on all
data up to time $t$, used to label the subsequent year, then refitted with the
new year appended. We require a 504-day (two-year) minimum training window and refit
every 252 days. This guarantees that all regime labels used in the backtest are strictly
out-of-sample---the single most important safeguard against look-ahead bias.

\paragraph{(e) Number of states.}
We fix the number of regimes at three on both theoretical and empirical grounds. Three
states map cleanly onto the Calm/Transitional/Crisis taxonomy, and the information
criteria in Figure~\ref{fig:modelsel} corroborate the choice: the Bayesian and Akaike
information criteria for both the HMM and the GMM flatten around $k=3$, while the GMM
silhouette score peaks there, indicating little gain from additional states relative to
their parameter cost.

\begin{figure}[H]
\centering
\includegraphics[width=0.49\textwidth]{hmm_bic_aic}\hfill
\includegraphics[width=0.49\textwidth]{gmm_bic_aic_silhouette}
\caption[Model selection criteria]{Model selection. Left: HMM BIC/AIC versus number of
states. Right: GMM BIC/AIC and silhouette scores.}
\label{fig:modelsel}
\floatnote{Source: own illustration based on own calculations.}
\end{figure}

\subsection{Regime-conditional CAPM}
To characterise the economic content of the regimes, we estimate, for each sector $i$ and
each VIX-defined regime $s$, the conditional Capital Asset Pricing Model (CAPM) by
ordinary least squares (OLS),
\begin{equation}
  e_{i,t}=\alpha_i^{(s)}+\beta_i^{(s)}\,e_{\mathrm{SPY},t}+\varepsilon_{i,t},
\end{equation}
reporting annualised alphas, betas, $t$-statistics and $R^2$.

\subsection{Portfolio construction}
Each strategy is rebalanced monthly (every 21 trading days) using two long-only
optimisers, solved by sequential least-squares quadratic programming (SLSQP):
\begin{align}
  \text{MVP:}\quad & \min_{w}\; w^{\top}\Sigma w
    && \text{s.t.}\;\; \mathbf{1}^{\top}w=1,\;\; w\ge 0, \\[2pt]
  \text{TPF:}\quad & \max_{w}\; \frac{\mu^{\top}w-\rf}{\sqrt{w^{\top}\Sigma w}}
    && \text{s.t.}\;\; \mathbf{1}^{\top}w=1,\;\; w\ge 0 .
\end{align}
The MVP depends only on $\Sigma$; the TPF additionally depends on the expected-return
vector $\mu$. Covariance matrices are estimated with Ledoit--Wolf linear shrinkage and
floored to positive-definiteness; both optimisers fall back to $1/N$ if the numerical
solve fails. For the regime-conditional strategies, $\mu$ and $\Sigma$ are estimated from
the historical days sharing the \emph{current} regime (30-observation minimum, else the
full history), so that all three regime classifiers feed an identical portfolio
construction and differ only in how they label the prevailing regime. During Crisis
regimes, 15\% of capital
is moved to the risk-free asset. This yields ten strategies: passive S\&P~500
buy-and-hold; the $1/N$ equal-weight portfolio; unconditional (``Static'') MVP and TPF;
and regime-conditional MVP and TPF for each of the three classifiers.

\subsection{Backtest, costs and metrics}
The backtest holds weights fixed between monthly rebalances and charges one-way
transaction costs proportional to realised turnover,
\begin{equation}
  \tau_t=\tfrac{1}{2}\sum_i\bigl|w_{i,\text{new}}-w_{i,\text{old}}\bigr|,
\end{equation}
at a baseline of 10 basis points; both net and gross daily return series are retained. We
report annualised return and volatility, the Sharpe and Sortino ratios, maximum drawdown,
Jensen's $\alpha$ and market $\beta$, the Calmar ratio, and the break-even transaction
cost---the total one-way cost (from gross returns) at which a strategy's Sharpe ratio
falls to that of the S\&P~500.

To guard against data snooping we compute, for the best strategy, the \textbf{Deflated
Sharpe Ratio} \citep{bailey2014} and \textbf{White's Reality Check} \citep{white2000}, a
stationary-bootstrap test \citep{politis1994} of the null that the best strategy does not
out-return the benchmark once the full search is accounted for. With daily
relative-performance series $d_{k,t}=r_{k,t}-r_{\mathrm{SPY},t}$, the Reality Check
statistic is
\begin{equation}
  V=\max_{k}\;\sqrt{T}\,\bar d_{k},
  \qquad
  V^{*}_{b}=\max_{k}\;\sqrt{T}\,\bigl(\bar d^{*}_{k,b}-\bar d_{k}\bigr),
  \qquad
  p=\frac{1}{B}\sum_{b=1}^{B}\mathbf{1}\{V^{*}_{b}\ge V\},
\end{equation}
using a mean block length of 10 days and $B=2{,}000$ resamples.

% =====================================================================
\section{Results}\label{sec:results}
% =====================================================================

\subsection{Regime characterisation}
The VIX rule assigns 68.6\% of days to Calm (mean VIX 14.8), 23.3\% to Transitional
(23.9), and 8.1\% to Crisis (40.7). The regimes are economically coherent
(Table~\ref{tab:regstats}): the equal-weight sector return averages $+27\%$ annualised in
Calm, $-9\%$ in Transitional and $-90\%$ in Crisis, while volatility rises monotonically.
Figure~\ref{fig:timeline} maps the three classifiers onto the VIX path; the crisis
clusters coincide with the 2008--09 global financial crisis and the early-2020 pandemic
shock. Figure~\ref{fig:heatmap} shows that the classifiers agree most in extreme states
and diverge in the Transitional middle ground, while Figure~\ref{fig:retdist} confirms
that the regime labels separate the cross-section of daily returns into progressively
wider and more left-skewed distributions.

\begin{table}[H]
\centering
\caption{VIX-defined regime statistics (full sample).}
\label{tab:regstats}
\small
\begin{tabular}{lrrrrr}
\toprule
Regime & Days & \% Days & Mean VIX & EqWt Ret. & EqWt Vol. \\
\midrule
Calm         & 3{,}795 & 68.6 & 14.8 & $+27.0\%$ & 10.2\% \\
Transitional & 1{,}292 & 23.3 & 23.9 & $-8.8\%$  & 19.3\% \\
Crisis       &   447   &  8.1 & 40.7 & $-90.0\%$ & 45.6\% \\
\bottomrule
\end{tabular}
\floatnote{Notes. Days are classified by the VIX rule (Calm $<20$, Transitional
$20$--$30$, Crisis $\ge 30$). ``EqWt Ret.'' and ``EqWt Vol.'' are the annualised mean and
standard deviation of the equal-weight sector return within each regime. Source: own
calculations based on data from Yahoo!\ Finance and FRED.}
\end{table}

The HMM recovers a more balanced and highly persistent partition (Calm 2{,}590 days,
Transitional 1{,}838, Crisis 601). Its sorted transition matrix (Table~\ref{tab:trans})
shows three features: states are persistent (diagonal $0.81$--$0.90$); the economy never
jumps \emph{directly} between Calm and Crisis but must pass through Transitional in both
directions---the ``ramp-up / ramp-down'' dynamic a memoryless classifier cannot
capture---and the Gaussian-mixture (GMM) labels (Calm 3{,}078,
Transitional 1{,}513, Crisis 438) lack this temporal structure.

\begin{table}[H]
\centering
\caption{Sorted HMM transition matrix.}
\label{tab:trans}
\small
\begin{tabular}{lccc}
\toprule
From $\backslash$ To & Calm & Transitional & Crisis \\
\midrule
Calm         & 0.905 & 0.095 & 0.000 \\
Transitional & 0.160 & 0.810 & 0.030 \\
Crisis       & 0.000 & 0.158 & 0.842 \\
\bottomrule
\end{tabular}
\floatnote{Notes. Entries are estimated one-day transition probabilities
$A_{ij}=\Pr(s_{t+1}=j\mid s_t=i)$ from the three-state Gaussian HMM, with states sorted by
their conditional mean $\dvix$. Rows sum to one. Source: own calculations.}
\end{table}

\begin{figure}[H]
\centering
\includegraphics[width=\textwidth]{01_regime_timeline}
\caption[Regime timeline]{Regime timeline: VIX level with the three classifiers' regime
assignments, 2004--2025.}
\label{fig:timeline}
\floatnote{Source: own illustration based on data from Yahoo!\ Finance and FRED.}
\end{figure}

\begin{figure}[H]
\centering
\includegraphics[width=0.86\textwidth]{02_regime_comparison_heatmap}
\caption[Classifier agreement]{Cross-classifier agreement on daily regime labels (VIX,
HMM, ML).}
\label{fig:heatmap}
\floatnote{Notes. Cells report the share of days on which two classifiers assign the same
regime label. Source: own illustration.}
\end{figure}

\begin{figure}[H]
\centering
\includegraphics[width=0.92\textwidth]{04_return_distributions}
\caption[Return distributions by regime]{Distribution of daily sector returns conditional
on the identified regime.}
\label{fig:retdist}
\floatnote{Source: own illustration based on data from Yahoo!\ Finance.}
\end{figure}

Diagnostic checks confirm the adequacy of the fitted models.
Figure~\ref{fig:hmmdiag} examines the HMM: the implied stationary (ergodic) distribution
matches the realised regime frequencies closely, and the autocorrelations of the
standardised residuals are small, indicating that the three-state Gaussian specification
captures the bulk of the serial structure in the feature process.

\begin{figure}[H]
\centering
\includegraphics[width=0.49\textwidth]{hmm_persistence_ergodic}\hfill
\includegraphics[width=0.49\textwidth]{hmm_residual_diagnostics}
\caption[HMM diagnostics]{HMM diagnostics. Left: regime persistence and ergodic
(stationary) distribution. Right: autocorrelation of standardised residuals.}
\label{fig:hmmdiag}
\floatnote{Source: own illustration based on own calculations.}
\end{figure}

\subsection{Regime-conditional betas}
The conditional CAPM (Table~\ref{tab:betas}, Figure~\ref{fig:betas}) reproduces the
canonical crisis-correlation pattern. Cyclical sectors' market betas rise into
crises---Financials from $1.04$ (Calm) to $1.41$ (Crisis) and Energy from $0.98$ to
$1.27$---so that diversification deteriorates precisely when protection is most valuable
\citep{ang2002}. Defensive sectors remain low-beta throughout (Utilities $\approx0.51$,
Consumer Staples $\approx0.55$, Health Care $\approx0.69$), and regression $R^2$ rises
markedly in crises (e.g.\ Technology $0.73\to0.91$), reflecting the collapse of
idiosyncratic dispersion under stress. This is the economic mechanism a successful
defensive rotation must exploit.

\begin{table}[H]
\centering
\caption{Regime-conditional CAPM betas (representative sectors).}
\label{tab:betas}
\small
\begin{tabular}{lccc}
\toprule
Sector & $\beta^{\text{Calm}}$ & $\beta^{\text{Trans}}$ & $\beta^{\text{Crisis}}$ \\
\midrule
Technology (XLK) & 1.18 & 1.13 & 1.00 \\
Financials (XLF) & 1.04 & 1.25 & 1.41 \\
Energy (XLE)     & 0.98 & 1.03 & 1.27 \\
\midrule
Health Care (XLV)     & 0.77 & 0.69 & 0.69 \\
Consumer Staples (XLP)& 0.55 & 0.53 & 0.57 \\
Utilities (XLU)       & 0.51 & 0.53 & 0.75 \\
\bottomrule
\end{tabular}
\floatnote{Notes. Each $\beta^{(s)}$ is the slope of an OLS regression of the sector's
daily excess return on the S\&P~500 excess return, estimated on the subset of days in
VIX regime $s$. The upper block shows cyclical sectors, the lower block defensive
sectors. Source: own calculations.}
\end{table}

\begin{figure}[H]
\centering
\includegraphics[width=0.92\textwidth]{05_capm_betas}
\caption[Regime-conditional betas]{Regime-conditional CAPM betas by sector: cyclical
betas rise into crisis, defensive betas stay low.}
\label{fig:betas}
\floatnote{Source: own illustration based on own calculations.}
\end{figure}

\subsection{Portfolio performance}
Table~\ref{tab:perf} reports out-of-sample performance over the common backtest window
(January~2006--December~2025, 5{,}029 days), net of 10~bps one-way costs.

\begin{table}[H]
\centering
\caption{Out-of-sample portfolio performance, net of transaction costs.}
\label{tab:perf}
\footnotesize
\setlength{\tabcolsep}{4.5pt}
\resizebox{\textwidth}{!}{%
\begin{tabular}{lrrrrrrrrr}
\toprule
Strategy & Ann.\ Ret.\ \% & Vol.\ \% & Sharpe & Sortino & Max DD \% & $\alpha$ \% & $p(\alpha)$ & Turn.\ \% & Break-even \\
\midrule
SPY B\&H    & 7.00 & 19.78 & 0.44 & 0.41 & $-59.6$ & 0.00 & --    & 0   & --    \\
EW $1/N$    & 5.71 & 19.08 & 0.39 & 0.36 & $-58.5$ & $-0.84$ & 0.38 & 0 & --    \\
Static MVP  & 6.11 & 14.53 & 0.48 & 0.45 & $-38.2$ & 1.52 & 0.38  & 76  & 86.1  \\
Static TPF  & 5.97 & 19.85 & 0.39 & 0.36 & $-48.8$ & 0.93 & 0.74  & 325 & 0.0   \\
VIX-MVP     & 6.04 & 13.44 & 0.50 & 0.47 & $-34.4$ & 1.88 & 0.28  & 73  & 125.3 \\
VIX-TPF     & 6.05 & 16.03 & 0.45 & 0.41 & $-32.2$ & 1.45 & 0.50  & 328 & 12.8  \\
HMM-MVP     & 6.19 & 13.95 & 0.50 & 0.47 & $-38.7$ & 2.03 & 0.29 & 122 & 77.5 \\
HMM-TPF     & 5.16 & 15.73 & 0.40 & 0.37 & $-55.1$ & 0.48 & 0.81 & 464 & 0.0 \\
\textbf{GMM-MVP} & \textbf{6.69} & 13.91 & \textbf{0.54} & 0.50 & $-39.8$ & \textbf{2.40} & 0.19 & 121 & \textbf{118.0} \\
GMM-TPF     & 3.86 & 18.14 & 0.30 & 0.28 & $-44.5$ & $-0.93$ & 0.71 & 433 & 0.0 \\
\bottomrule
\end{tabular}}
\floatnote{Notes. Daily out-of-sample excess returns, net of 10~bps one-way transaction
costs, over 2006--2025 ($T=5{,}029$). Sharpe and Sortino are annualised; $\alpha$ and
$p(\alpha)$ are the annualised Jensen alpha and its two-sided $p$-value from a CAPM
regression on the S\&P~500. ``Turn.'' is measured annual one-way turnover; ``Break-even''
is the total one-way cost (bps), computed from gross returns, at which the Sharpe ratio
equals the S\&P~500's. Source: own calculations.}
\end{table}

Four patterns emerge. \emph{First}, the value of regime conditioning is overwhelmingly
defensive: every minimum-variance strategy compresses volatility to roughly $13$--$15\%$
(versus $20\%$ for the S\&P~500) and the maximum drawdown to about $-34\%$ to $-40\%$
(versus $-60\%$), while raising the Sharpe ratio above the benchmark's $0.44$ to between
$0.48$ and $0.54$ (Figures~\ref{fig:wealth}--\ref{fig:dd}). \emph{Second}, na\"ive
diversification disappoints: the $1/N$ portfolio underperforms the cap-weighted S\&P~500 on
every risk-adjusted metric, a caution against reading \citet{demiguel2009} as an
unconditional endorsement of equal weighting. \emph{Third}, and most tellingly, because all
three classifiers feed an identical portfolio construction, the tangency (maximum-Sharpe)
variants fail \emph{uniformly}: every TPF strategy underperforms both its own
minimum-variance counterpart and the S\&P~500's $7\%$ return, with Sharpe ratios of only
$0.30$--$0.45$ (the GMM tangency portfolio is weakest, at $0.30$). Estimating expected
returns from the regime-conditional historical mean is simply too noisy to time, whichever
classifier supplies the regime. \emph{Fourth}, the minimum-variance strategies, which never
touch the fragile expected-return estimate, are uniformly strong and remarkably similar
across classifiers: GMM-MVP leads (Sharpe $0.54$, annual alpha $2.40\%$), with HMM-MVP and
VIX-MVP both at $0.50$. That the \emph{memoryless} GMM edges out the Markovian HMM, and that
even the parameter-free VIX rule matches it, indicates that the gains come from \emph{risk
targeting under any reasonable regime map}, not from the sophistication of the regime model.
Figure~\ref{fig:weights} traces how the regime-conditional weights rotate toward defensive
sectors and cash as the system enters Crisis.

\begin{figure}[H]
\centering
\includegraphics[width=\textwidth]{07_cumulative_wealth}
\caption[Cumulative wealth]{Cumulative wealth of the benchmark and the regime-conditional
strategies over the backtest window.}
\label{fig:wealth}
\floatnote{Source: own illustration based on own calculations.}
\end{figure}

\begin{figure}[H]
\centering
\includegraphics[width=\textwidth]{09_drawdown_profiles}
\caption[Drawdown profiles]{Drawdown profiles: the minimum-variance strategies roughly
halve the peak drawdown of the passive benchmark.}
\label{fig:dd}
\floatnote{Source: own illustration based on own calculations.}
\end{figure}

\begin{figure}[H]
\centering
\includegraphics[width=\textwidth]{06_portfolio_weights}
\caption[Portfolio weight evolution]{Evolution of regime-conditional portfolio weights
across the backtest.}
\label{fig:weights}
\floatnote{Notes. Stacked sector weights for the regime-gated strategies; the cash sleeve
appears during Crisis regimes. Source: own illustration.}
\end{figure}

\subsection{Transaction costs and break-even}
Measuring turnover directly is essential, because the strategies differ by an order of
magnitude (Table~\ref{tab:perf}). The minimum-variance strategies trade modestly
($73$--$122\%$ per year), but the tangency portfolios churn aggressively---between $325\%$
and $464\%$ of capital annually. This asymmetry compounds the tangency portfolios' weak
gross performance: the break-even analysis (Figure~\ref{fig:breakeven}) shows that every
TPF strategy fails to match the S\&P~500 Sharpe \emph{even at zero cost} (break-even
$\le 13$~bps), whereas the minimum-variance strategies are highly cost-tolerant---VIX-MVP
absorbs $125$~bps of one-way cost, GMM-MVP $118$~bps, and Static MVP $86$~bps before their
Sharpe falls to the benchmark's. The cost-robust strategies are precisely the low-turnover
minimum-variance ones.

\begin{figure}[H]
\centering
\includegraphics[width=0.92\textwidth]{10_breakeven_costs}
\caption[Break-even transaction cost]{Break-even transaction-cost analysis: Sharpe ratio
(from gross returns) versus total one-way cost.}
\label{fig:breakeven}
\floatnote{Notes. The dashed line marks the S\&P~500 Sharpe ratio; a strategy's break-even
cost is the abscissa at which its curve crosses it. Source: own illustration.}
\end{figure}

\subsection{Statistical significance}
No strategy's CAPM alpha is significant at the 5\% level on a single-equation basis; the
best-performing strategy, GMM-MVP, comes closest at $p=0.19$. The relevant question,
however, is whether the \emph{best of all the strategies searched} is significant, which
requires a multiple-testing correction (Table~\ref{tab:sig}).

\begin{table}[H]
\centering
\caption{Data-snooping corrections for the best strategy (GMM-MVP).}
\label{tab:sig}
\small
\begin{tabular}{lll}
\toprule
Test & Statistic & Verdict \\
\midrule
Deflated Sharpe Ratio & $0.97$ (Sharpe $0.54$ vs.\ null $0.11$) & Marginally significant at 95\% \\
White's Reality Check & $p=0.98$ (2{,}000 bootstraps) & Cannot reject $H_0$ \\
\bottomrule
\end{tabular}
\floatnote{Notes. The Deflated Sharpe Ratio \citep{bailey2014} reports the probability
that the true Sharpe exceeds the snooping-adjusted null benchmark, correcting for the
number of trials and the skewness and kurtosis of returns. White's Reality Check
\citep{white2000} reports the stationary-bootstrap $p$-value (mean block 10 days,
$B=2{,}000$) of the null that the best strategy does not out-return the S\&P~500.
Source: own calculations.}
\end{table}

The two corrections paint a coherent picture. The Deflated Sharpe Ratio of the
highest-Sharpe strategy, GMM-MVP, is $0.97$: its annualised Sharpe of $0.54$ clears the
snooping-adjusted null threshold of $0.11$, so its \emph{risk-adjusted} performance is
marginally significant even after accounting for the breadth of the search. White's Reality
Check, however, is unambiguous: with a bootstrap $p$-value of $0.98$, we \emph{cannot}
reject the null that the best strategy fails to out-\emph{return} the S\&P~500, reflecting
that no strategy in the study beats the benchmark on raw return at all. Together the two
tests locate the value firmly in the \emph{denominator}: the best strategy earns a return
comparable to the index with materially less risk---a statistically defensible improvement
in risk-adjusted terms, but no evidence of return outperformance once data snooping and
serial dependence are accounted for.

% =====================================================================
\section{Discussion}\label{sec:disc}
% =====================================================================
The results cohere around a single economic message. The conditional-beta evidence shows
that the exploitable structure in the volatility cycle is the rising market sensitivity of
cyclical sectors under stress; the portfolio evidence shows that the strategies which
monetise this structure are the low-turnover, risk-targeting ones; and the significance
evidence shows that what survives scrutiny is risk reduction, not return generation. This
is consistent with a market in which volatility and correlation are far more forecastable
than the sign of returns---a stylised fact underpinning the practical success of
minimum-variance investing \citep{clarke2006} and the difficulty of beating simple
benchmarks with return forecasts \citep{demiguel2009}.

The comparison across classifiers is instructive. The HMM's transition matrix encodes
genuinely useful structure---the obligatory passage through the Transitional state---yet
this temporal information buys remarkably little. For risk targeting, the memoryless GMM
(Sharpe $0.54$) actually edges out the Markovian HMM ($0.50$), with the parameter-free VIX
rule ($0.50$) level---so modelling persistence is, if anything, a slight handicap here
rather than a driver. For return timing, the classifier matters in a different way---every
tangency variant fails (Sharpe $0.30$--$0.45$, all below the benchmark's return),
confirming that the binding constraint is the noisiness of regime-conditional expected
returns, not the quality of the regime labels. The robust, defensible gains live entirely in the minimum-variance column,
and are available from the crude VIX rule almost as readily as from the HMM or the GMM. For
a practitioner, the policy implication is to prefer a simple, low-turnover, regime-aware
minimum-variance overlay over any return-timing apparatus.

% =====================================================================
\section{Limitations and Future Work}\label{sec:limit}
% =====================================================================
Several caveats temper the conclusions. (i) The backtest holds weights fixed between
monthly rebalances and measures turnover on a target-to-target basis, abstracting from
intra-month drift; realistic implementation turnover would be somewhat higher,
reinforcing the case against the high-turnover variants. (ii) Transaction costs are a
constant 10~bps proportional charge and exclude market impact. (iii) The crisis cash
overlay and the $20/30$ VIX thresholds are fixed ex ante. (iv) The Deflated Sharpe Ratio
treats the strategies searched as the trial universe, but the true degrees of freedom of
the research process are larger. (v) The HMM is estimated with full covariance matrices and
roughly 200 parameters on initial windows of 504 observations, a regime in which
covariance estimates are noisy. Promising directions include soft-information portfolios
that weight by regime \emph{posterior probabilities}, a Hidden semi-Markov specification
with explicit duration distributions, and an out-of-sample combination of the three
classifiers.

% =====================================================================
\section{Conclusion}\label{sec:concl}
% =====================================================================
Conditioning a US sector-rotation strategy on inferred market regimes---whether by a simple
VIX rule, a Gaussian Hidden Markov Model, or an unsupervised Gaussian Mixture
Model---delivers a large, robust, and economically meaningful reduction in portfolio risk:
drawdowns fall from roughly $-60\%$ to about $-37\%$ and Sharpe ratios rise from $0.44$ to
$0.48$--$0.54$, at low and cost-tolerant turnover. The same cannot be said for return
enhancement. Because every classifier feeds an identical regime-conditional portfolio
construction, the return-timing (tangency) variants fail uniformly---none beats the
benchmark's return---and no strategy out-returns the S\&P~500 once data snooping is
corrected for. Two conclusions follow. First, the value of regime conditioning resides in
the choice of portfolio objective---minimum-variance risk targeting---rather than in return
timing. Second, that value is strikingly insensitive to the sophistication of the regime
model: the memoryless mixture matches or slightly exceeds the Markov-switching model, and
even the fixed VIX rule is level with it. In this dataset, the intelligence in ``intelligent'' regime
allocation is best spent on managing risk.

% =====================================================================
\clearpage
\addcontentsline{toc}{section}{References}
\begin{thebibliography}{99}
% =====================================================================
\bibitem[Ang and Bekaert(2002)]{ang2002} Ang, A., and Bekaert, G. (2002). International Asset Allocation with Regime Shifts. \textit{Review of Financial Studies}, 15(4), 1137--1187.
\bibitem[Bailey and L\'opez de Prado(2014)]{bailey2014} Bailey, D.~H., and L\'opez de Prado, M. (2014). The Deflated Sharpe Ratio: Correcting for Selection Bias, Backtest Overfitting, and Non-Normality. \textit{Journal of Portfolio Management}, 40(5), 94--107.
\bibitem[Baum et al.(1970)]{baum1970} Baum, L.~E., Petrie, T., Soules, G., and Weiss, N. (1970). A Maximization Technique Occurring in the Statistical Analysis of Probabilistic Functions of Markov Chains. \textit{Annals of Mathematical Statistics}, 41(1), 164--171.
\bibitem[Bloom(2009)]{bloom2009} Bloom, N. (2009). The Impact of Uncertainty Shocks. \textit{Econometrica}, 77(3), 623--685.
\bibitem[Clarke et al.(2006)]{clarke2006} Clarke, R., de Silva, H., and Thorley, S. (2006). Minimum-Variance Portfolios in the US Equity Market. \textit{Journal of Portfolio Management}, 33(1), 10--24.
\bibitem[DeMiguel et al.(2009)]{demiguel2009} DeMiguel, V., Garlappi, L., and Uppal, R. (2009). Optimal Versus Naive Diversification: How Inefficient Is the 1/N Portfolio Strategy? \textit{Review of Financial Studies}, 22(5), 1915--1953.
\bibitem[Dempster et al.(1977)]{dempster1977} Dempster, A.~P., Laird, N.~M., and Rubin, D.~B. (1977). Maximum Likelihood from Incomplete Data via the EM Algorithm. \textit{Journal of the Royal Statistical Society, Series B}, 39(1), 1--38.
\bibitem[Guidolin and Timmermann(2007)]{guidolin2007} Guidolin, M., and Timmermann, A. (2007). Asset Allocation under Multivariate Regime Switching. \textit{Journal of Economic Dynamics and Control}, 31(11), 3503--3544.
\bibitem[Hamilton(1989)]{hamilton1989} Hamilton, J.~D. (1989). A New Approach to the Economic Analysis of Nonstationary Time Series and the Business Cycle. \textit{Econometrica}, 57(2), 357--384.
\bibitem[Harvey et al.(2016)]{harvey2016} Harvey, C.~R., Liu, Y., and Zhu, H. (2016). \ldots and the Cross-Section of Expected Returns. \textit{Review of Financial Studies}, 29(1), 5--68.
\bibitem[Ledoit and Wolf(2004)]{ledoit2004} Ledoit, O., and Wolf, M. (2004). A Well-Conditioned Estimator for Large-Dimensional Covariance Matrices. \textit{Journal of Multivariate Analysis}, 88(2), 365--411.
\bibitem[Markowitz(1952)]{markowitz1952} Markowitz, H. (1952). Portfolio Selection. \textit{Journal of Finance}, 7(1), 77--91.
\bibitem[McLachlan and Peel(2000)]{mclachlan2000} McLachlan, G., and Peel, D. (2000). \textit{Finite Mixture Models}. New York: Wiley.
\bibitem[Politis and Romano(1994)]{politis1994} Politis, D.~N., and Romano, J.~P. (1994). The Stationary Bootstrap. \textit{Journal of the American Statistical Association}, 89(428), 1303--1313.
\bibitem[Sharpe(1966)]{sharpe1966} Sharpe, W.~F. (1966). Mutual Fund Performance. \textit{Journal of Business}, 39(1), 119--138.
\bibitem[Viterbi(1967)]{viterbi1967} Viterbi, A.~J. (1967). Error Bounds for Convolutional Codes and an Asymptotically Optimum Decoding Algorithm. \textit{IEEE Transactions on Information Theory}, 13(2), 260--269.
\bibitem[Whaley(2009)]{whaley2009} Whaley, R.~E. (2009). Understanding the VIX. \textit{Journal of Portfolio Management}, 35(3), 98--105.
\bibitem[White(2000)]{white2000} White, H. (2000). A Reality Check for Data Snooping. \textit{Econometrica}, 68(5), 1097--1126.
\end{thebibliography}

% =====================================================================
\clearpage
\addcontentsline{toc}{section}{Statement of Academic Integrity}
\section*{Statement of Academic Integrity}
% =====================================================================
Hereby, I declare that I have composed the presented paper independently on my own and
without any other resources than the ones indicated. All thoughts taken directly or
indirectly from external sources are properly denoted as such. This paper has neither been
previously submitted to another authority nor has it been published yet.

\vspace{0.8em}
\textit{Declaration on the use of generative AI.} In the preparation of this paper, the
author used the generative AI system Claude (Anthropic) for the following purposes:
assistance in structuring and debugging the Python implementation; drafting and editing
text passages that were subsequently reviewed and revised by the author; and assistance in
producing tables and figures from the author's own data and computations. All AI-assisted
outputs were critically reviewed, verified against the underlying data and source code, and
edited by the author, who takes full responsibility for the content of this paper.
\emph{[Adapt this paragraph to reflect your actual use; remove if not applicable.]}

\vspace{2.0cm}
\noindent\begin{tabular}{p{6cm}p{6cm}}
\hrulefill & \hrulefill \\
Place, Date & Signature \\
\end{tabular}

\end{document}
